\scanpage{93}
Ces dernières formules sont surtout intéressantes dans le contexte où on travaille \deleted{\unsure{la|ici}} sur \deleted{les}la base absolue $\widetilde S_0$ et avec les structures de Legendre\inserted{-Jacobi} directes, et où le schéma modulaire $U_{03}\times_{S_0}\mu_4^\times$ sur $S_0$ est remplacé par $U_{03\,\widetilde S_0}$ $(=U_{03}\times_{S_0}\mu_4^\times\,!)$ sur $\widetilde S_0$, et le groupe modulaire $\mathbb D_{3\,S_0}\times_{S_0}\mu_{4\,S_0}$ est remplacé par le \inserted{ss-}groupe $\mathbb D'_{3\,\widetilde S_0}$ de \deleted{\unsure{\mathbb D_3|\mathbb D_8}} son image inverse sur $\widetilde S_0$.

\begin{center}--\quad --\quad --\end{center}

Il faut maintenant définir les épinglages de Jacobi d'échelon $4$ relatifs sur l'échelon $2$, \deleted{i.e. les sections} \inserted{qui sont des sections \inserted{(\unsure{variables|véritables})}} de $ {}_4E\longrightarrow {}_2E$ au dessus de $\unsure{\underline I_E|\underline J_E}=({}_2E)^*$, \unsure{satisfont|satisfant|satisfaisant} la condition exprimée \deleted{\unsure{que|par}} sur chaque fibre géométrique par $\unsure{\zeta_i+\zeta_j+\zeta_k|\zeta_i+\zeta_j+\zeta_h|\xi_i+\xi_j+\xi_k}=0$, et montrer comment un tel épinglage définit canoniquement un épinglage de Legendre.
% archive sheet number: 92
